%% Методические указания
%% Лаборатория математики ФН1
\documentclass[12pt, a4paper]{article}
\usepackage[T2A]{fontenc}
\usepackage[utf8]{inputenc}
\usepackage[russian]{babel}
\usepackage[left=25mm, right=20mm, top=20mm, bottom=20mm]{geometry}
\usepackage{amsmath, amssymb}
\usepackage{enumitem}
\pagestyle{plain}
\setlength{\parindent}{1.25cm}
\begin{document}
\begin{center}
  {\large\bfseries Методические указания}

  \vspace{0.3em}
  {\large Работа с графами: обход и характеристики}

  \vspace{0.5em}
  {\small Дискретная математика · 1 курс\\ Лаборатория математики ФН1}
\end{center}
\vspace{1em}

\section*{Что проверять в задаче о графе}

\begin{enumerate}[nosep]
  \item \textbf{Степени вершин.} Сумма степеней равна удвоенному числу рёбер
        (лемма о рукопожатиях)~--- это первая проверка правильности рисунка.
  \item \textbf{Связность.} Обходом в ширину или глубину найдите компоненты связности.
  \item \textbf{Эйлеров цикл} существует тогда и только тогда, когда граф связен
        и все степени чётны. Эйлеров путь~--- если ровно две вершины нечётной степени.
  \item \textbf{Гамильтонов цикл} простого критерия не имеет. Пользуйтесь достаточными
        условиями (теорема Дирака: все степени $\ge n/2$) либо перебором с обоснованием.
  \item \textbf{Остовное дерево} стройте поиском в ширину от указанной вершины,
        помечая порядок обхода.
\end{enumerate}

\section*{Как оформлять}

\begin{itemize}[nosep]
  \item Вершины подписывайте теми же именами, что в матрице смежности.
  \item Рисунок делайте без пересечений рёбер, если это возможно.
  \item Для обхода приводите порядок посещения вершин, а не только результат.
  \item Вывод о существовании цикла обосновывайте ссылкой на критерий,
        а не словами «видно из рисунка».
\end{itemize}

\section*{Типичные ошибки}

\begin{itemize}[nosep]
  \item Петля даёт вклад 2 в степень вершины, а не 1.
  \item Из отсутствия эйлерова цикла заключают об отсутствии гамильтонова. Это разные свойства.
  \item Остовное дерево строят «на глаз», не указывая алгоритм.
\end{itemize}

\end{document}
